\section{The Ladder}

\textbf{[B]} The lepton masses can be written as rungs of a $\xi$-ladder:
\[
  m_i = r_i\,\xi^{p_i}\,v ,
\]
with rational structure numbers $r_i$, rational exponents $p_i$, and a bridge
constant $v$. For the leptons these are
\[
  \bigl(r_e,p_e\bigr) = \Bigl(\tfrac43,\tfrac32\Bigr),\quad
  \bigl(r_\mu,p_\mu\bigr) = \Bigl(\tfrac{16}{5},1\Bigr),\quad
  \bigl(r_\tau,p_\tau\bigr) = \Bigl(\tfrac{25}{9},\tfrac23\Bigr).
\]

\textbf{[K]} In the \emph{ratios} $v$ cancels, leaving parameter-free expressions:
\[
  \frac{m_\mu}{m_e} = \frac{12}{5}\,\xi^{-1/2},\qquad
  \frac{m_\tau}{m_\mu} = \frac{125}{144}\,\xi^{-1/3},\qquad
  \frac{m_\tau}{m_e} = \frac{25}{12}\,\xi^{-5/6}.
\]

According to A080 these are the intrinsic quantities: they depend on the geometry
and nothing else. They are precisely what is to be measured against.

\section{Where the Exponents Come From}

\textbf{[K]} The exponents $p_i$ are not free. They are winding numbers on a
\emph{logarithmic} torus with step size $\Delta p = 1/3$. This step size is the
fundamental third of the harmonic series -- the same geometry of periodic boundary
conditions that generates overtones on a vibrating string. Half-integer values
are geometric means between two winding levels, negative values counter-windings.

\textbf{[K]} The three lepton exponents lie on this ladder:
\[
  p_e = \tfrac32,\qquad p_\mu = 1,\qquad p_\tau = \tfrac23,
\]
with steps
\[
  p_e - p_\mu = \tfrac12 \ (\text{geometric mean}),\qquad
  p_\mu - p_\tau = \tfrac13 \ (\text{fundamental third}).
\]
These are the same winding numbers $(n_\theta, n_\phi)$ that also fix the prefactors
$r_i = n_\phi^2/n_\theta$ (A110): mass prefactor and exponent come from one source,
not two. The ladder is therefore not a fit but the reading of the resonance
structure of the torus.

\section{What the Ladder Delivers}

\textbf{[K]} The calculation against PDG values gives:

\begin{center}
\begin{tabular}{lrrr}
\hline
Ratio & Ladder & PDG & Deviation \\
\hline
$m_\mu/m_e$    & $207{.}85$ & $206{.}77$ & $+0{.}52\,\%$ \\
$m_\tau/m_\mu$ & $16{.}992$ & $16{.}817$ & $+1{.}04\,\%$ \\
$m_\tau/m_e$   & $3532$     & $3477$     & $+1{.}57\,\%$ \\
\hline
\end{tabular}
\end{center}

\textbf{[K]} The tolerance standard from A010 and A080 requires three figures
side by side. The measurement precision of the comparison quantities is at
$m_\mu/m_e$ in the range $10^{-9}$, at $m_\tau$ about $7\cdot10^{-5}$. The
determination tolerance of the ladder values is orders of magnitude coarser: the
structure numbers are declared rational, not determined to significant figures.

\textbf{Assessment:} the ladder carries the \emph{hierarchy} and the per cent
level. It does not carry precision. A one per cent deviation is not a finding
against the structure here, because the structure claims nothing more precise at
this level -- but it is also not a success that could be reported to five figures.

\section{The Reverse Direction: What the Data Return for $\xi$}

\textbf{[K]} Solving the three ratios for $\xi$, the data give
\[
  \xi_{\mu/e} = 1{.}347\cdot10^{-4},\quad
  \xi_{\tau/\mu} = 1{.}375\cdot10^{-4},\quad
  \xi_{\tau/e} = 1{.}358\cdot10^{-4},
\]
against $4/30000 = 1{.}3333\cdot10^{-4}$, i.e.\ $+1{.}1$ to $+3{.}2\,\%$.
The three values scatter among themselves by about $2\,\%$.

This is the honest empirical status: the ladder is consistent in the ratios at
the per cent level, not exactly. It does not return a uniform $\xi$.

\textbf{[K]} The same from the other side: solving for the bridge constant per
lepton gives $248{.}9$, $247{.}6$, and $245{.}1$ GeV, scattered around
$246{.}22$ GeV. The residual deviation is therefore absorbed into the bridge
constant and does \emph{not} cancel completely in the ratios.

\section{The Countercheck That Fails}

\textbf{[K]} The obvious conjecture, that the residual deviations follow the
fractal correction factor, is \emph{denied} by the calculation. $K_\text{frak}
= 1-100\xi$ corresponds to $-1{.}333\,\%$, the half and three-half powers to
$-0{.}665$ and $-1{.}995\,\%$. The measured deviations are $+0{.}52$ and
$+1{.}04\,\%$ -- they fit neither in magnitude nor in \emph{sign} a single
$K_\text{frak}$ power.

This section is here because a failed countercheck belongs to the finding just
as much as a successful one. The obvious explanatory attempt has been tested and
has failed; it is not rescued by a free exponent.

\textbf{The striking doubling is arithmetically forced -- and does not evidence a
correction.} The $\tau/\mu$ deviation is almost exactly double the $\mu/e$
deviation ($1{.}04\,\%$ vs.\ $0{.}52\,\%$), and the $\tau/e$ deviation is their
sum ($0{.}52 + 1{.}04 = 1{.}56\,\%$). This is not a new law and not a
coincidence but a pure bookkeeping necessity: along a multiplicative ladder in
which $m_\tau/m_e = (m_\tau/m_\mu)(m_\mu/m_e)$ holds, the logarithmic deviations
add. Of the three deviations therefore only \emph{one} is independently measured
(the ratio $1{.}98$); the third is fixed as a sum.

\textbf{Importantly, this step does \emph{not}} establish anything about the
fractal factor that the countercheck just rejected -- on the contrary, it shows
that the doubling singles out no correction at all: it already follows from the
ladder form alone. One who reads a law from the doubling reads a single measured
number three times.

\textbf{[K]} The factor that carries the ladder per generation (its
generation-counting step, order $N_0 \approx 39$) is a \emph{different} quantity
from the fixed scale $100$ in $K_\text{frak}$: the one is the ladder step, the
other the RG scale of the absolute correction (A040). The value of this ladder
step is not derived -- that is the one open point named below.

\section{Two Mutually Exclusive Computation Modes}

\textbf{[STIPULATION]} Before any of these deviations is booked, the computation
mode must be fixed. There are two, and they exclude each other.

\textbf{Mode 1, reference-free.} All ratios are computed purely from $\xi$,
without anchor. Then all three stand independently against the data, each
$0{.}5$ to $1{.}6\,\%$ off, and all three are legitimate \emph{check points}.

A note on terminology, because otherwise a contradiction with A010 arises: these
three ratios are \emph{not} prediction quantities of the series. The prediction
list comprises four quantities, and the lepton masses appear there solely via the
$\tau$-mass from the circulant (A110), not via the ladder. The ladder provides
check points at the per cent level -- less than a prediction and more than nothing.

\textbf{Mode 2, anchored.} $m_\mu/m_e$ is declared a reference point. Then this
ratio is \emph{spent}: its deviation may no longer be booked as a test result,
otherwise the same quantity is used twice -- once as anchor, once as test.

\textbf{[B]} Mixing the two modes is the most common error in this sector. The
A-series consistently uses mode 1, unless something else is explicitly stated.

\section{Verification Script}

\begin{itemize}
  \item All numbers in this document -- ladder ratios, reverse-$\xi$, bridge
    constants per lepton, the failed $K_\text{frak}$ countercheck:
    \texttt{python/A\_Serie\_Skripte/a100\_leiter.py}
\end{itemize}

\section{Symbol Glossary}

Symbols used throughout the entire A-series are explained in A010.
Specific to this document:

\begin{center}
\renewcommand{\arraystretch}{1.3}
{\small\begin{tabular}{lp{9cm}}
\toprule
Symbol & Meaning \\
\midrule
$r_i\in\mathbb{Q}$ & Rational prefactor in $m_i=r_i\cdot\xi^{p_i}\cdot v$ \\
$p_i\in\mathbb{Q}$ & Exponent; ladder step size $\Delta p=1/3$ \\
$(n,l,j)$ & Principal, orbital, total angular momentum of the mode \\
$v=246\,\text{GeV}$ & Higgs vacuum expectation value as energy unit \\
\bottomrule
\end{tabular}}
\renewcommand{\arraystretch}{1.0}
\end{center}

\section{Open Questions}

\textbf{First, the per cent scatter is intended, not a defect.} The ladder is a
statement at the per cent level; a $\xi$ back-computed from it scatters accordingly
by $2\,\%$. For this reason the sector anchors the absolute scale not in the
ladder but in a single reference point (A110): the ladder carries the hierarchy,
the reference point the precision.

\textbf{Second, the exponents $p_i$ are derived, not merely motivated.} They are
winding numbers on the logarithmic torus (above), with step size $1/3$ and
geometric means at half-integer values -- the same source from which the prefactors
$r_i$ follow. The condition from A020 is thereby met: the exponent is fixed
independently, the $\xi$ power is a statement.

\textbf{Third, the doubling is clarified.} It is arithmetically forced (above):
along a multiplicative ladder the logarithmic deviations add, only one is
independent. No unbooked residual remains.

\textbf{Fourth, exactly one point remains open -- and it is named.} The one
independently measured deviation corresponds to a generation-counting ladder step
of order $N_0 \approx 39$, whose value is not derived. It is to be separated
from the fixed RG scale $100$ in $K_\text{frak}$ (A040) and is not explained by
the bridge constant in which the residual deviation is absorbed (A080). The value
of this step and the candidates for it are collected in A105; whether structure
sits behind it is the only genuinely open question of the sector.

\section{Supersedes}

This document subsumes: Dok.~006 (mass formula), Dok.~046 (mass calculation),
Dok.~016 (particle masses), Dok.~189 (winding numbers of exponents), Dok.~190
(R55 additivity, P42 reference point), and the ladder sections from Dok.~292.
Incorporated: K2 (exponent $p_\tau$), P40 (level separation), P42 (reference
point), P35, R55 (evidential breadth: only one point independently measured).

\section{Use of AI Tools}

This document was created with the assistance of AI tools. The questions,
stipulations, and all substantive decisions come from the author; the tools
formulate, compute, and create verification scripts.
Responsibility for content and errors rests entirely with the author.
Full wording of the rules in A010.
